% P16-C4: graph-uniform quantifiers and the modified Tian algorithm.
\section{Graph-uniform identification and the modified Tian algorithm}
\label{sec:p16-graph-uniform-identification}

The four-valued positive hedge witnesses already settle a complete
graph-level identification question. They do so even while an additional
theorem for some prescribed binary or ternary domains remains separate.
This section states that distinction and supplies the modified Tian
algorithm in Shpitser and Pearl's Figure~5, including its termination,
soundness and completeness proofs
\cite[Figure~5, Theorems~7--8 and Corollary~5,
pp.~1225--1226]{ShpitserPearl2006JointIdentification}.
The source permits nonpositive observed laws and adds an action-support
restriction. The strictly positive execution domain used here is an
explicit handbook hypothesis, not a hypothesis attributed to that paper.

\begin{hbdefinition}{def:p16-graph-uniform-positive}{Graph-uniform positive-law identification}
Fix a finite observed graph \(G\) on \(N\) and disjoint action
and response sets \(X,Y\subseteq N\). For each family
\(\mathbf S=(S_i)_{i\in N}\) of finite nonempty observed value sets,
use the graph-defined model class \(\mathcal M^+(G,\mathbf S)\)
from \cref{def:p16-hedge-identification-class}. Latent value sets,
independent root distributions and total equations may vary within
that class; no extra latent-incidence or cardinality restriction is imposed.

The query is \emph{graph-uniformly identified on positive laws} if
\[
 \begin{gathered}
 \text{for every such }\mathbf S\text{ and every }
 M_1,M_2\in\mathcal M^+(G,\mathbf S),\\
 p_{M_1}=p_{M_2}
 \quad\Longrightarrow\quad
 P^{M_1}_x(y)=P^{M_2}_x(y)
 \text{ for every }x\in S_X, y\in S_Y.
 \end{gathered}
\]
Models compared with each other use the same \(\mathbf S\).
The outer quantifier requires this property for every finite alphabet
choice. Fixed-family identification omits that outer quantifier.
Equivalently, for every \(\mathbf S\) the compatible positive
observed laws determine a unique intervention kernel through a
family-indexed function \(\Phi_{G,X,Y,\mathbf S}\).

A \emph{uniform finite formula} is a finite expression in observed
masses, products, ratios and sums over the named value sets, whose
graphical construction is independent of the alphabet sizes. Its
sums are evaluated on the particular input \(\mathbf S\).
A uniform do-calculus derivation uses such expressions and the same
finite graph-based rule licenses for every \(\mathbf S\).
\end{hbdefinition}

\begin{hbtheorem}{thm:p16-graph-uniform-equivalences}{proved=annex:p16.graph-uniform-equivalences}{The graph-uniform ID, hedge and do-calculus criteria}
For the fixed \(G,X,Y\), the following are equivalent:
\begin{enumerate}
\item The query is graph-uniformly identified on positive laws.
\item ID succeeds on the graph and these action/response sets.
\item There is no hedge in \(G\) for this complete query \((X,Y)\).
\item A uniform finite do-calculus derivation with ordinary probability
      operations identifies the kernel from the observed law.
\end{enumerate}
An ID failure is a graph decision: its occurrence does not depend
on the particular compatible positive table or on the finite alphabet
sizes. Every failure yields an original-query hedge and two positive
four-valued models with the same observations and different effects.
Consequently it disproves the outer universal quantifier already;
no fixed-binary or fixed-ternary extension is required for this result.

For a \emph{specified} alphabet family containing singleton domains,
the earlier singleton elimination can strengthen the answer for that
family. Graph-uniform identification retains the original graph because
its outer quantifier also includes families without those constants.
\end{hbtheorem}

\begin{hbcorollary}{thm:p16-graph-uniform-all-effects}{proved=annex:p16.graph-uniform-all-effects}{Uniform identification of every unconditional effect}
Every disjoint action/response query in \(G\) is graph-uniformly
identified on positive laws if and only if no directed edge has both
endpoints in the same maximal district of \(G\). Equivalently,
there is no vertex connected to one of its directed children by a
bidirected path. When this condition holds, ID and do-calculus
identify every query by uniform finite formulas.
\end{hbcorollary}

\begin{hbtheorem}{thm:p16-graph-uniform-action-scope}{proved=annex:p16.graph-uniform-action-scope}{Why the hedge criterion retains the entire action set}
In the graph \(A\to M\to Y\), \(A\leftrightarrow M\),
the kernel \(P_a(Y)\) is not graph-uniformly identified on
positive laws, whereas
\[
 P_{a,m}(y)=p(y\mid m)
\]
is graph-uniformly identified. Thus a hedge for a smaller action
set need not obstruct a larger intervention.

For an explicit positive witness, let \(U,E,E_Y\) be independent
bits, with \(U\) fair and \(E,E_Y\) of one-probability \(1/4\).
Model 1 uses \(A=U\), \(M=A\oplus U\oplus E\),
\(Y=M\oplus E_Y\); model 2 replaces \(M\) by \(E\).
Their common observed masses, for each \(a\), are
\[
\begin{array}{c|cc}
m&p(a,m,0)&p(a,m,1)\\\hline
0&9/32&3/32\\
1&1/32&3/32
\end{array}
\]
but \(P^1_a(Y=1)=1/2\) and \(P^2_a(Y=1)=3/8\).
Under \(\operatorname{do}(A=a,M=m)\), both give response-one
probability \(1/4\) at \(m=0\) and \(3/4\) at \(m=1\).
\end{hbtheorem}

\subsection{The modified Tian identification procedure}

\begin{hbdefinition}{def:p16-tian-kernel-interface}{Current Q kernels and fixed arguments}
Fix an original positive observed law \(p\), graph \(G\) on
\(N\), value family \(\mathbf S\), and one topological order.
All functions below are functions of the full assignment \(v\in S_N\).
For \(T\subseteq N\), the kernel \(q_T(v_T;v_{N\setminus T})\)
has random coordinates \(T\) and fixed intervention parameters
\(N\setminus T\). Its invariant is
\[
 q_T(v)=Q[T](v)>0\qquad\text{for every full assignment }v
\]
in every compatible original model. It is normalized over \(S_T\)
at each fixed outside assignment. These outside arguments are not
observational conditioning events.

Use \(\Sigma_B\) for the explicit coordinate-sum operator in
\cref{def:p16-id-call-data}. For a maximal district \(D\) of
\(G[T]\), write \(\operatorname{Ext}^T_D(q_T)\) for the
proved product of prefix conditional ratios from
\cref{thm:p16-latent-district-recovery}, using the restricted
topological order and retaining every outside parameter.

A component call has nonempty bidirected-connected sets
\(C\subseteq T\), the induced graph \(G[T]\), and this
kernel. Its target is \(Q[C]\), with every coordinate in
\(N\setminus C\) a fixed parameter. A failure record retains
the local pair \(G[T],G[C]\), the current kernel and its fixed
set, all containing call frames, and the original query and district
selected by the outer procedure when there is one.
\end{hbdefinition}

\begin{hbmethod}{def:p16-tian-component-algorithm}{The component procedure c-identify}
On input \((C,T,q_T)\) with the preceding interface, put
\(A=\operatorname{an}_{G[T]}(C)\), and apply these branches in order:
\begin{enumerate}
\item If \(A=C\), return \(\Sigma_{T\setminus C}q_T\).
      This includes \(C=T\).
\item If \(A=T\), return the failure record
      \(\mathsf{FAIL}(G[T],G[C])\).
\item Otherwise \(C\subsetneq A\subsetneq T\). Form
      \(q_A=\Sigma_{T\setminus A}q_T\). Let \(T'\) be the
      unique maximal district of \(G[A]\) containing \(C\),
      and form \(q_{T'}=\operatorname{Ext}^A_{T'}(q_A)\).
      Return \(\mathsf{c\mbox{-}identify}(C,T',q_{T'})\).
\end{enumerate}
In the recursive branch the exact inclusions are
\[
 C\subseteq T'\subseteq A\subsetneq T.
\]
The selected district may equal \(C\) or \(A\). A failed
child propagates its record with this parent frame; it is never
treated as an expression to be summed or divided.
\end{hbmethod}

\begin{hblemma}{thm:p16-tian-component-invariant}{proved=annex:p16.tian-component-invariant}{The component recursion is defined and preserves positive Q kernels}
Every branch is defined. The recursive \(T'\) exists uniquely,
is nonempty and bidirected-connected, contains \(C\), and is
strictly smaller than the previous \(T\). Its kernel is exactly
\(Q[T']\), positive and normalized for every retained outside
assignment. All conditional denominators used in extraction are
positive. At failure, \(C\subsetneq T=\operatorname{an}_{G[T]}(C)\).
\end{hblemma}

\begin{hbtheorem}{thm:p16-tian-component-soundness}{proved=annex:p16.tian-component-soundness}{Component termination and success soundness}
Every component call terminates, with depth at most
\(|T|-|C|+1\). Every successful return is exactly \(Q[C]\)
as a function of the full assignment, including its outside parameters.
Its expression uses finite sums, products and positive-denominator ratios
of the input kernel. Its steps have derivations through the already
proved ancestral and district do-calculus identities.
\end{hbtheorem}

\begin{hbmethod}{def:p16-tian-outer-algorithm}{The outer procedure identify}
Input consists of \(G,N,X,Y,\mathbf S,p\), with disjoint
\(X,Y\subseteq N\) and compatible \(p>0\) everywhere.
\begin{enumerate}
\item If \(Y=\varnothing\), return one. Otherwise put
      \(D=\operatorname{an}_{G[N\setminus X]}(Y)\).
\item Compute the maximal districts \(D_1,\ldots,D_k\) of
      \(G[D]\) and the maximal districts of \(G\). For each
      \(D_i\), let \(T_i\) be the unique full-graph district
      containing it, and compute
      \(q_{T_i}=\operatorname{Ext}^N_{T_i}(p)\).
\item Obtain \(e_i=\mathsf{c\mbox{-}identify}(D_i,T_i,q_{T_i})\).
      If a child fails, return its record together with the original
      query and selected \(D_i\). Otherwise return
      \[
       e(v)=\left(\Sigma_{D\setminus Y}\prod_{i=1}^k e_i\right)(v).
      \]
\end{enumerate}
The action argument is \(v_X=x\), the response argument is
\(v_Y=y\), and summation binds exactly \(D\setminus Y\).
Values on \(N\setminus(D\cup X)\) remain displayed parameters
until the final expression; the next theorem proves independence
of them, so they may then be set to any reference tuple.
\end{hbmethod}

\begin{hbtheorem}{thm:p16-tian-outer-soundness}{proved=annex:p16.tian-outer-soundness}{Outer termination, target soundness and support}
The outer procedure terminates. Every successful return equals
\(P_x(y)\) for every compatible model on the specified value sets
and every \(x,y\). It is independent of the other displayed
parameters. All reached kernels satisfy the stated positive support
invariant, and the successful trace gives a uniform finite expression
and do-calculus derivation. A specified singleton-domain family may
first use the proved constant-coordinate reduction.
\end{hbtheorem}

\begin{hbtheorem}{thm:p16-tian-failure-hedge}{proved=annex:p16.tian-failure-hedge}{A component failure gives an original-query hedge}
Suppose a component called for the outer district \(C=D_i\)
fails at \(T\). Then \(T\cap X\ne\varnothing\), while
\(C\cap X=\varnothing\). Let \(K\) be the directed sink
set of \(G[C]\). Choose directed edges within \(C\) toward
\(K\), and from every vertex of \(T\setminus C\) toward
\(C\). Extend a bidirected spanning tree of \(C\) to one
on \(T\). These selected edges give \(K\)-rooted forests
on \(T,C\) which form a hedge for the original \((X,Y)\).
No equality between the unpruned sink sets of \(G[C]\) and
\(G[T]\) is assumed.

For a standalone component call with initial random set \(T_0\),
the same failure gives a hedge in \(G[T_0]\) for action
\(T_0\setminus C\) and response \(C\).
\end{hbtheorem}

\begin{hbtheorem}{thm:p16-tian-completeness}{proved=annex:p16.tian-completeness}{Graph-uniform completeness of identify and c-identify}
The outer procedure succeeds exactly for the graph-uniformly identified
positive-law queries. It therefore agrees with ID on success versus
failure for every graph and action/response pair, and its successful
expression gives the same intervention kernel. These expressions
need not have identical syntax. In particular, identify succeeds for
every query exactly when no directed edge lies within one maximal district.

The standalone component procedure is complete for recovering the
complementary-intervention kernel on \(C\) from the positive law
on \(T_0\), graph-uniformly over finite nonempty alphabets.
Here \(G[T_0]\) is the entire observed graph of the standalone
problem; any outside parameters have been fixed in advance.
For this standalone question, failure also proves nonidentification
for every prescribed family on \(T_0\) whose domains all have
at least two values: the extracted roots already belong to the
response \(C\), so the positive binary witness and alphabet lift
suffice. These are model-class statements about the available input
law, without an assumption that arbitrary externally specified kernels
can be embedded into a particular additional model.
\end{hbtheorem}

\subsection{Worked branches and source-domain boundary}

\begin{hbtheorem}{thm:p16-tian-worked}{proved=annex:p16.tian-worked}{Both equality cases, a front-door return and a root-set distinction}
The recursive district can realize both allowed equalities.
\begin{enumerate}
\item With directed edge \(X\to Y\) and bidirected path
      \(X\leftrightarrow Z\leftrightarrow Y\), the component
      call \(C=\{Y\},T=\{X,Z,Y\}\) has
      \(A=\{X,Y\}\) and \(T'=C\). A positive example has
      \(p(x,z,y)=5/32\) when \(z=y\) and \(3/32\) otherwise.
      The call returns \(Q[\{Y\}]=1/2\).
\item With directed edge \(X\to Y\) and bidirected path
      \(X\leftrightarrow Y\leftrightarrow Z\), the same component
      input has \(A=\{X,Y\}\) and \(T'=A\). With common
      positive masses \(p(x,0,z)=3/16\), \(p(x,1,z)=1/16\),
      the recursive call fails on \(T'=\{X,Y\},C=\{Y\}\).
\item For \(A\to B\to C,A\leftrightarrow C\), identifying
      \(P_a(c)\) returns
      \[
       \sum_b p(b\mid a)\sum_{a'}p(a')p(c\mid a',b).
      \]
      On the positive model of \cref{thm:p16-id-worked-success},
      its response-one values are \(19/32\) and \(25/32\).
\item Let the directed graph be \(a\to t\to b\), with
      \(a\leftrightarrow b\) and \(a\leftrightarrow t\).
      For component \(C=\{a,b\}\) and \(T=\{a,t,b\}\),
      one has \(T=\operatorname{an}_{G[T]}(C)\), but the
      unpruned sink sets are \(\{a,b\}\) and \(\{b\}\).
      The common-root construction deletes \(a\to t\), keeps
      \(t\to b\), and has roots \(\{a,b\}\) in both forests.
\end{enumerate}
\end{hbtheorem}

\begin{proofboundary}{def:p16-graph-uniform-source-regime}{The source's nonpositive and action-support regime}
The source explicitly permits zero observed masses and writes an
additional requirement
\(P(x\mid\operatorname{pa}_G(X)\setminus X)>0\)
\cite[p.~1220 and Lemma~1, p.~1221]{ShpitserPearl2006JointIdentification}.
Every conditional expression also requires a positive conditioning mass.
These are distinct from the full positivity assumed by the execution
theorems above. The parent-set expression must be interpreted at its
parent assignments; for example, a precise stratum-wise admissibility
condition would require
\[
 p(x,b)>0\quad\text{for every }b\in S_{\operatorname{pa}_G(X)\setminus X}
                         \text{ with }p(b)>0.
\]
This original-action condition does not license every recursive
extraction: \cref{thm:p16-zero-front-door-obstruction} supplies two
models satisfying it with equal observations and different effects.
\Cref{thm:p16-zero-support-quantifiers} distinguishes the two
parent-assignment readings and the stronger quantification over all actions.

The supported extension is supplied by
\cref{thm:p16-zero-masked-extraction,thm:p16-zero-algorithm-soundness}:
each current complementary extraction uses its own proved recovery
conditions, and only a complete recovered kernel is passed onward.
The licensed nonpositive trace domain and its graph-uniform
completeness theorem are \cref{thm:p16-zero-licensed-completeness}.
All equation-replacement interventions remain structurally defined,
including at observationally absent assignments. Identification of every
individual mass at every arbitrary sparse law is a further question.
\end{proofboundary}
